\chapter{Further Reading and References}
\markboth{Further Reading and References}{Further Reading and References}
\begingroup
\raggedbottom

The following books develop complementary parts of the framework used here. Choose a source for the particular claim or calculation being investigated, and consult its assumptions, notation, edition, and errata. A robotics or control theorem does not become evidence about human golf performance merely because it can be applied to a simplified swing model.

\section*{Geometry, Kinematics, and Mechanical Models}

\begin{itemize}
\item \textbf{A Mathematical Introduction to Robotic Manipulation}, Murray, Li, and Sastry \cite{Murray1994}, treats rigid-body motion, manipulator and hand kinematics, dynamics, and nonholonomic motion planning. The \href{https://www.cds.caltech.edu/~murray/mlswiki/index.php?title=Main_Page}{authors' companion site} provides chapter information and corrections. Use it to connect Lie groups and Jacobians to explicitly defined frames and mechanical systems.
\item \textbf{Geometric Control of Mechanical Systems}, Bullo and Lewis \cite{BulloLewis2004}, develops differential-geometric modeling and control of mechanical systems. Its \href{https://motion.me.ucsb.edu/book-gcms/}{companion page} includes contents, selected material, examples, and errata. This is a suitable route into connections, mechanical control structures, and the hypotheses behind geometric analysis.
\item \textbf{Robot Modeling and Control}, Spong, Hutchinson, and Vidyasagar \cite{Spong2006}, covers kinematics, Jacobians, trajectory planning, dynamics, joint and multivariable control, force control, and vision-based methods. The \href{https://bcs.wiley.com/he-bcs/Books?action=contents&bcsId=2888&itemId=0471649902}{publisher's chapter listing} identifies the scope of the cited edition. Its robot models still require adaptation and validation before supporting a biomechanical interpretation.
\end{itemize}

\section*{Computational Rigid-Body Dynamics}

\begin{itemize}
\item \textbf{Rigid Body Dynamics Algorithms}, Featherstone \cite{Featherstone2008}, develops spatial-vector conventions and recursive dynamics algorithms. The \href{https://royfeatherstone.org/teaching/index.html}{author's teaching materials} and \href{https://royfeatherstone.org/}{software resources} support implementation. Distinguish the linear-cost articulated-body calculation for a tree from the additional work needed for closed loops, constraints, and contacts.
\end{itemize}

\begin{samepage}
\section*{Nonlinear Stability and Feedback}

\begin{itemize}
\item \textbf{Applied Nonlinear Control}, Slotine and Li \cite{SlotineLi1991}, develops nonlinear stability and feedback methods. \href{https://www.mit.edu/~nsl/videos/lectures.html}{Slotine's accompanying course materials} identify the relevant book sections and separate the later contraction lectures and papers. Check the domain and model conditions of a stability argument before extending it to saturated or hybrid systems.
\item \textbf{Contraction Theory for Dynamical Systems}, Bullo \cite{BulloContraction}, develops contractivity and incremental stability using specified norms and metrics. The \href{https://fbullo.github.io/ctds}{author's book page} provides version information and materials. The bibliography cites the 2024 edition; later revisions may reorganize or extend the treatment. Contraction claims must retain their metric, input, domain, and existence assumptions.
\end{itemize}

The chapter-level primary references supply more specialized results for transverse dynamics, phase constraints, nonlinear funnels, stochastic models, and learning. Keep mathematical derivations, numerical verification, model calibration, and experimental validation identifiable when combining those sources.

\nocite{LohmillerSlotine1998, ManchesterSlotine2017, Khatib1987, Shiriaev2010}

\end{samepage}
\clearpage
\endgroup
